\section{Stage 2 results: does the direction steer?}
\label{sec:stage2results}

This is the pivotal section. With the operating point frozen on dev and git-tagged, we touch the
held-out set once and evaluate the four-conjunct criterion. The held-out arms---unsteered-\NOSYS,
steered-\NOSYS, norm-matched random ($k{=}5$), and an HF \FULL{} reference---are compared on both
oracle signals in Figure~\ref{fig:F8}; steering reproduces $\numslot{STEER-REPRO}$ of the skill's
POC gain. The result resolves into one of three preregistered branches, written out in full below.
Exactly one is kept on real data.

\subsection{Branch A --- a causally sufficient (and/or necessary) design direction exists}
\label{sec:branchA}

\begin{resultbranch}{A: causally sufficient direction}
Steering the diff-in-means direction with \emph{no design prompt} raises held-out UI quality on both
oracle signals: steered-\NOSYS{} beats unsteered-\NOSYS{} by $\numslot{STEER-DELTA}$ on the POC
(significant on $\numslot{STEER-NFAM}$ of six objective families after Holm) and is preferred with BT
utility $\numslot{STEER-BT}$ (cluster-bootstrap CI $>0$). Crucially, the norm-matched random control
does not clear both bars ($\numslot{STEER-RAND}$ POC, CI bracketing zero), so the effect is a
\emph{direction} effect, not a \emph{norm} effect---a random push of equal magnitude does not make
pages better. The direction reproduces $\numslot{STEER-REPRO}$ of the full skill's held-out gain
(Figure~\ref{fig:F8}), and it does so on prompts and \emph{task types the extraction corpus never
contained}. To our knowledge this is the first demonstration of a causal, steerable direction for the
visual/aesthetic quality of generated front-end code: a few sentences of design guidance and a
single mid-network vector turn out to move the same underlying quantity, and that quantity can be
actuated with the prose removed.

\smallskip\noindent\emph{Coherence and specificity.} The gain is not bought with gibberish: at
$\rhostar$ the mean next-token KL from the unsteered model is $\numslot{STEER-KL}$ nats (under the
$0.30$ guardrail) and steered render-success is $\numslot{STEER-RENDER}$. Steering leaves general
ability intact---non-UI algorithmic pass-rate drops by only $\numslot{SPEC-PASSDROP}$ (within the
$10\%$ tolerance) and general-text KL is $\numslot{SPEC-KL}$---so the direction is specific to design,
not a blunt instrument that degrades everything.
\end{resultbranch}

\paragraph{Necessity (the flip test).} Projecting $\vhat$ out at every layer and position while the
skill is present attenuates the skill's own gain by $\numslot{FLIP-ATTEN}$ (Eq.~\ref{eq:flip}), which
promotes the claim from sufficiency to necessity.
\begin{resultbranch}{A.1: sufficient and necessary}
The flip test attenuates the skill's gain by $\numslot{FLIP-ATTEN}$ with CI $>0$: removing the single
direction from the residual stream substantially disables the skill even when its 391 tokens are still
in context. The direction is therefore both sufficient (adding it with no prompt helps) and necessary
(removing it stops the prompt from helping)---the strong two-arm result, mirroring the refusal-direction
template but for a graded, multi-component aesthetic behavior. This is the paper's strongest mechanistic
claim: the skill's effect is, to a large extent, \emph{routed through this one direction}.
\end{resultbranch}
\begin{resultbranch}{A.2: sufficient, not demonstrably necessary}
Addition succeeds but the flip test does not attenuate the skill's gain ($\numslot{FLIP-ATTEN}$, CI
includes zero): the direction is a sufficient carrier but not the unique one. The skill routes its
effect through this direction \emph{and} others---consistent with a concept cone rather than a single
axis. We report ``sufficient, not demonstrably necessary'' and do not overclaim uniqueness; the honest
framing is ``\emph{a} causally sufficient design direction,'' never ``\emph{the} design direction.''
\end{resultbranch}

\paragraph{Failure surface and generalization (RQ7).} A positive mean can hide heavy-tailed per-prompt
behavior, so we report the distribution. The overall anti-steerable fraction (prompts whose steered POC
falls below unsteered) is $\numslot{FAIL-ANTI-OVERALL}$; on unseen types it is
$\numslot{FAIL-ANTI-UNSEEN}$, and the mean unseen-type steering delta is $\numslot{FAIL-UNSEEN-D}$
(Figure~\ref{fig:F12}). The reproduction fraction on unseen types is $\numslot{FAIL-REPRO-UNSEEN}$.
\begin{resultbranch}{RQ7 robust}
The anti-steerable fraction stays below $50\%$ across all task types, seen and unseen, and the
unseen-type delta is positive and gated: the direction transfers cleanly out-of-distribution. Steering a
direction extracted only from marketing- and content-oriented dev pages nonetheless improves
control-dense settings panels and admin tables, evidence that it encodes transferable design
\emph{content} rather than a task-idiom artifact.
\end{resultbranch}
\begin{resultbranch}{RQ7 partial}
Steering helps on average and generalizes to most unseen types, but a named regime is brittle:
$\numslot{FAIL-WORST-TYPE}$ crosses the $50\%$ anti-steerable line ($\numslot{FAIL-WORST-ANTI}$), so
steering makes those pages worse as often as better (Figure~\ref{fig:F12}). This is the honest failure
map---the direction is not a blanket win, and its data-dense failure mode is exactly where the aesthetic
prior (generous whitespace, one signature element) conflicts with the task's information density.
\end{resultbranch}
\begin{resultbranch}{RQ7 brittle}
The unseen-type effect collapses (anti-steerable $\ge50\%$ overall on OOD): the direction is
dev-specific and does not transport to genuinely new layouts. We scope the causal claim to the
in-distribution regime and report the OOD boundary precisely; the direction is real but local, and its
extraction corpus bounds where it works.
\end{resultbranch}

\paragraph{Component$\leftrightarrow$direction correspondence (RQ8).} Finally we test whether the
component structure the ablation found survives into the representation. Component sub-vectors
$v_{C_i}=\bar h_{\lstar}^{(\AOI\text{-}C_i)}-\bar h_{\lstar}^{(\NEUTRAL)}$ have mean off-diagonal
cosine $\numslot{CORR-COSMEAN}$ (max $\numslot{CORR-COSMAX}$), and $\numslot{CORR-NMATCH}$ of five steer
their own metric family (Figure~\ref{fig:F9}). Cosine can show sub-vectors are linearly distinguishable
but not that each \emph{causally} steers its family, so the rule pairs near-orthogonality with a
steering-signature match.
\begin{resultbranch}{RQ8 found}
The sub-vectors are near-orthogonal (mean $|\!\cos|=\numslot{CORR-COSMEAN}<0.3$) \emph{and}
$\ge3/5$ steer their own metric family---C1 moves color, C2 moves layout, and so on. The design skill
has a \emph{component basis} inside the model, and the black-box ablation signatures correspond to
white-box steering directions: the paper's strongest bridge result, an ablation$\leftrightarrow$activation
correspondence predicted by the geometry of linear representations. It means ``clean design'' is not one
entangled thing but a small set of separable, individually actuatable sub-concepts.
\end{resultbranch}
\begin{resultbranch}{RQ8 partial}
One or two sub-vectors steer their own family while the rest are entangled ($\numslot{CORR-NMATCH}/5$
match): the representation has partial component structure. Some components (likely color and layout,
the most necessary) have distinct directions; others share a subspace. The bridge holds in part---the
most causal components are the most separable---but ``clean design'' is not a fully diagonal basis.
\end{resultbranch}
\begin{resultbranch}{RQ8 absent}
The sub-vectors are not near-orthogonal, or fewer than three steer their own family: there is no
component basis. ``Clean design'' is one entangled direction inside the model even though the prompt-level
skill decomposes into five components. This is still informative and slightly surprising: the ablation's
component structure is a property of the \emph{instruction}, not of the model's internal
representation---the model compresses five instructions into one axis of variation.
\end{resultbranch}

\subsection{Branch B --- the honest null}
\label{sec:branchB}

\begin{resultbranch}{B: not linearly steerable}
No single linear direction steers held-out UI quality under the coherence guardrail. Steered-\NOSYS{}
does not beat unsteered-\NOSYS{} on any objective family after Holm ($\numslot{STEER-DELTA}$ POC, CI
$\numslot{STEER-DELTA-CI}$ bracketing zero) and the preference delta's BT CI includes zero
(Figure~\ref{fig:F8}). This is a \emph{null with teeth}, not an absence of evidence: the power gate
certifies the confidence intervals are tight enough to exclude effects the study was designed to
detect, and the extraction direction \emph{is} decodable (RQ5 probes separate \FULL{} from \NEUTRAL{}
at high AUC)---so we have a direction that \emph{correlates} with design but does not \emph{cause} it
when injected. A correlational direction that does not steer is a hypothesis, not a finding; we report
it as such.

\smallskip\noindent The evidential value is precise. The null \emph{rules out} the strongest form of
the linear-representation account for this behavior in this model: that a single mid-network direction,
added with no prompt, suffices to reproduce the skill's effect. It is consistent with two readings we
cannot distinguish here and state honestly: either the skill's effect is \emph{distributed} across
many directions and/or layers (a single injection is too blunt---the pre-stated multi-layer and
subspace fallbacks were tried and also failed), or the behavior is genuinely \emph{non-linear} in the
residual stream (diff-in-means, a linear probe of a nonlinear model, cannot capture it). Either way,
the practical implication is that the design skill's mechanism is not a simple steerable feature: the
391 tokens do more than trigger one latent switch. For the interpretability literature, it is a
negative datapoint on how far the single-direction picture---so clean for refusal---extends to graded,
compositional, generation-spanning aesthetic behaviors. We publish it with the same rigor as a positive
result, because a forced or buried null would corrupt exactly the oracle this project is built to
defend.
\end{resultbranch}

\subsection{Branch C --- partial outcomes}
\label{sec:branchC}
Between a clean causal direction and a clean null lie four partial results, each shipped with its
numbers and each carrying a distinct, bounded claim.

\begin{resultbranch}{C.1: steers the objective signal only}
Steered-\NOSYS{} beats unsteered on $\ge1$ objective family after Holm ($\numslot{STEER-DELTA}$), but
the preference delta's BT CI includes zero. The direction moves the \emph{measurable} correlates of
design---contrast, alignment, slop tells---without moving \emph{human-judged} quality. Given that
objective metrics cap at $\sim$half the variance in appeal, this is a meaningful boundary: we have
steered what the metrics see but not what a designer prefers, and by the null rule the aesthetic claim
is inadmissible. We report a ``metric-level'' steering effect and explicitly decline the quality claim.
\end{resultbranch}
\begin{resultbranch}{C.2: steers, but only by breaking coherence}
A configuration raises POC only above the KL guardrail ($\numslot{STEER-KL}$ nats $>0.30$) or below the
render floor: the direction ``improves'' pages by pushing the model off-distribution, and the gain
rides on artifacts rather than design. Under the preregistered guardrail this does not count as
success. We report the direction as \emph{coherence-breaking} and show the dose-response
(Figure~\ref{fig:F7}) where the apparent gain and the KL blow-up coincide---an honest account of a
non-result that a laxer guardrail would have mis-sold as a finding.
\end{resultbranch}
\begin{resultbranch}{C.3: anti-steerable majority}
The mean steering effect is near zero but hides a heavy-tailed distribution: $\numslot{FAIL-ANTI-OVERALL}$
of prompts are anti-steerable, so steering helps half and harms half (Figure~\ref{fig:F12}). The
``average effect'' is an artifact of cancellation, and reporting only the mean would be actively
misleading---the reason the protocol reports the full per-prompt distribution. The direction is not a
reliable design lever; whatever it encodes is prompt-contingent in a way a single global vector cannot
resolve.
\end{resultbranch}
\begin{resultbranch}{C.4: sufficient but not necessary}
Addition succeeds on both signals (Branch~A holds) but the flip test does not attenuate the skill's
gain ($\numslot{FLIP-ATTEN}$, CI includes zero): a steerable direction exists, yet the skill does not
\emph{depend} on it---the model has other routes to the same behavior. This is the concept-cone
outcome; we claim a sufficient direction and an honest account of non-uniqueness, and we present the
cosine/cone structure (Figure~\ref{fig:F9}) rather than a single-axis story.
\end{resultbranch}

\synthfig[0.47\linewidth]{F8}{F8_preview.pdf}{8}{The held-out causal test. POC by arm---unsteered,
steered (no design prompt), norm-matched random ($k{=}5$), and \FULL{} reference---with the
\%-skill-reproduced annotation. Success requires steered $>$ unsteered while random $\approx$ unsteered.
Planted arms illustrate the layout; the real bars decide RQ6.}

\synthfig[0.53\linewidth]{F12}{F12_preview.pdf}{8}{Anti-steerable fraction by task type (seen vs.\
unseen/OOD), with the $50\%$ brittleness line. A positive mean can coexist with a brittle unseen type;
the distribution, not the mean, is the RQ7 claim. Planted; real data replaces the bars.}

\synthfig[0.7\linewidth]{F9}{F9_preview.pdf}{9}{Component$\leftrightarrow$direction correspondence
(RQ8). Left: cosine matrix among the five component sub-vectors and $v_{\text{full}}$ (near-zero
off-diagonal $\Rightarrow$ orthogonal). Right: steer-signature $\times$ ablation-signature match (a
bright diagonal $\Rightarrow$ each component steers its own family). Planted diagonal; real data
decides found/partial/absent.}
